
\documentclass{article}
\usepackage{colortbl}
\usepackage{makecell}
\usepackage{multirow}
\usepackage{supertabular}

\begin{document}

\newcounter{utterance}

\centering \large Interaction Transcript for game `cladder', experiment `full\_v1.5\_default', episode 9125 with Qwen3.5{-}9B.
\vspace{24pt}

{ \footnotesize  \setcounter{utterance}{1}
\setlength{\tabcolsep}{0pt}
\begin{supertabular}{c@{$\;$}|p{.15\linewidth}@{}p{.15\linewidth}p{.15\linewidth}p{.15\linewidth}p{.15\linewidth}p{.15\linewidth}}
    \# & \multicolumn{2}{c}{Player} && \multicolumn{2}{c}{Game Master} \\
    \hline

    \theutterance \stepcounter{utterance}  
    & & & \multicolumn{4}{p{0.6\linewidth}}{
        \cellcolor[rgb]{0.9,0.9,0.9}{
            \makecell[{{p{\linewidth}}}]{
                \texttt{\tiny{[P1$\langle$GM]}}
                \texttt{You are an expert in causal inference. The following question is not a typical commonsense query, but rather a meticulously designed question created by a professor specializing in causal inference, intended to assess the students' mastery of the course content.} \\
\texttt{Imagine a self{-}contained, hypothetical world with only the following conditions, and without any unmentioned factors or causal relationships: Kraz has a direct effect on pexu and rukz. Pexu has a direct effect on hwax. Hwax has a direct effect on rukz. Kraz is unobserved. For those who are not pexu and are not hwax, the probability of rukz is 19\%. For those who are not pexu and are hwax, the probability of rukz is 29\%. For those who are pexu and are not hwax, the probability of rukz is 20\%. For those who are pexu and are hwax, the probability of rukz is 31\%. For those who are not pexu, the probability of hwax is 16\%. For those who are pexu, the probability of hwax is 90\%. For those who are pexu, would it be less likely to see rukz if the individual was not pexu?} \\
\texttt{Start your answer with 'Yes' or 'No', followed by additional reasoning or evidence to support your explanation.} \\
            }
        }
    }
     \\ \\

    \theutterance \stepcounter{utterance}  
    & \multicolumn{4}{p{0.6\linewidth}}{
        \cellcolor[rgb]{0.561,0.737,0.561}{
            \makecell[{{p{\linewidth}}}]{
                \texttt{\tiny{[P1$\rangle$GM]}}
                \texttt{No} \\
\\ 
\texttt{**Reasoning:**} \\
\\ 
\texttt{To answer this question, we must first clarify the logical structure of the query within the constraints of the hypothetical world. The question asks: "For those who are **pexu**, would it be less likely to see **rukz** if the individual was **not pexu**?"} \\
\\ 
\texttt{This phrasing presents a logical contradiction regarding the state of the individual. We are asked to evaluate the probability of an outcome (\$Rukz\$) for a specific group defined by a condition (\$Pexu = True\$), but then ask how that probability changes if we hypothetically change the defining condition of that group to its negation (\$Pexu = False\$).} \\
\\ 
\texttt{In causal inference and probability theory, we cannot simultaneously hold that an individual *is* \$Pexu\$ and *is not* \$Pexu\$. The question essentially asks for a comparison between two mutually exclusive counterfactual states for the same subject, or it is a trick question relying on the definition of the set.} \\
\\ 
\texttt{However, interpreting the question as a request to compare the **actual probability** of \$Rukz\$ for the group \$Pexu\$ against the **actual probability** of \$Rukz\$ for the group \$\textbackslash{}neg Pexu\$ (which is the only meaningful comparison available given the data):} \\
\\ 
\texttt{1.  **Calculate \$P(Rukz | Pexu)\$:**} \\
\texttt{    The problem provides conditional probabilities for \$Rukz\$ based on the joint states of \$Pexu\$ and \$Hwax\$. We must sum the probabilities for the two cases where \$Pexu\$ is true:} \\
\texttt{    *   Case A: \$Pexu = True, Hwax = False\$. Probability of \$Rukz = 20\textbackslash{}\%\$.} \\
\texttt{    *   Case B: \$Pexu = True, Hwax = True\$. Probability of \$Rukz = 31\textbackslash{}\%\$.} \\
\\ 
\texttt{    To find the total probability for \$Pexu\$, we need the distribution of \$Hwax\$ given \$Pexu\$. The prompt states: "For those who are pexu, the probability of hwax is 90\%."} \\
\texttt{    \$\$P(Rukz | Pexu) = P(Rukz | Pexu, \textbackslash{}neg Hwax) \textbackslash{}times P(\textbackslash{}neg Hwax | Pexu) + P(Rukz | Pexu, Hwax) \textbackslash{}times P(Hwax | Pexu)\$\$} \\
\texttt{    \$\$P(Rukz | Pexu) = (0.20 \textbackslash{}times 0.10) + (0.31 \textbackslash{}times 0.90)\$\$} \\
\texttt{    \$\$P(Rukz | Pexu) = 0.02 + 0.279 = 0.299 \textbackslash{}text} \\
            }
        }
    }
    & & \\ \\

    \theutterance \stepcounter{utterance}  
    & & & \multicolumn{2}{p{0.3\linewidth}}{
        \cellcolor[rgb]{0.9,0.9,0.9}{
            \makecell[{{p{\linewidth}}}]{
                \texttt{\tiny{[GM$|$GM]}}
                \texttt{No} \\
            }
        }
    }
    & & \\ \\

    \theutterance \stepcounter{utterance}  
    & & & \multicolumn{2}{p{0.3\linewidth}}{
        \cellcolor[rgb]{0.9,0.9,0.9}{
            \makecell[{{p{\linewidth}}}]{
                \texttt{\tiny{[GM$|$GM]}}
                \texttt{yes} \\
            }
        }
    }
    & & \\ \\

    \theutterance \stepcounter{utterance}  
    & & & \multicolumn{2}{p{0.3\linewidth}}{
        \cellcolor[rgb]{0.9,0.9,0.9}{
            \makecell[{{p{\linewidth}}}]{
                \texttt{\tiny{[GM$|$GM]}}
                \texttt{game\_result = LOSE} \\
            }
        }
    }
    & & \\ \\

\end{supertabular}
}

\end{document}
